% Shared preamble of the proof overview documents  <Folder>/<Name>_overview.tex.
% Each overview starts with
%     % Shared preamble of the proof overview documents  <Folder>/<Name>_overview.tex.
% Each overview starts with
%     % Shared preamble of the proof overview documents  <Folder>/<Name>_overview.tex.
% Each overview starts with
%     \input{../../scripts/overview_preamble.tex}
% (for ThinMatching/: \input{../scripts/overview_preamble.tex}) and is compiled from its own folder by
% scripts/build_overviews.py.
\documentclass[10pt,a4paper]{article}
\usepackage[T1]{fontenc}
\usepackage[utf8]{inputenc}
\usepackage[margin=2.2cm]{geometry}
\usepackage{amsmath,amssymb,amsthm}
\usepackage{array,booktabs,longtable}
\usepackage[dvipsnames]{xcolor}
\usepackage{enumitem}
\usepackage{tikz}
\usetikzlibrary{arrows.meta,positioning,shapes.multipart,fit,calc}
\usepackage[hidelinks]{hyperref}
\setlist{itemsep=2pt,topsep=4pt}
\renewcommand{\arraystretch}{1.2}

% Lean names, files and paper references
\newcommand{\lean}[1]{\textcolor{blue!50!black}{\texttt{\detokenize{#1}}}}
\newcommand{\file}[1]{\textcolor{OliveGreen!70!black}{\texttt{\detokenize{#1}}}}
\newcommand{\paperref}[1]{\textsf{#1}}
\newcommand{\hyp}[1]{\textcolor{BrickRed}{\textbf{#1}}}

% Boxes for the dependency diagrams
\tikzset{
  res/.style={draw, rounded corners, fill=blue!6, align=center, inner sep=4pt, font=\small},
  main/.style={res, fill=blue!18, very thick},
  cited/.style={res, fill=red!10, dashed},
  shared/.style={res, fill=gray!12},
  dep/.style={-{Stealth[length=2.2mm]}, thick},
}

\newcommand{\overviewtitle}[2]{%
  \begin{center}{\Large\bfseries #1}\\[4pt]{\small Proof overview of the Lean~4/Mathlib formalisation}\\[2pt]
  {\small #2}\end{center}}

% Standard paragraph describing the paper
\newcommand{\paperversion}{\noindent\textbf{Paper.} J.~Breitung, T.~Gervens, M.~Grohe.
  \emph{Graph Similarity Meets Optimal Transport}.\\ File \file{Graph_Similarity_Meets_Optimal_Transport.pdf} in the
  repository; statement, equation, page and line numbers refer to this PDF. \file{docs/NUMBERING.md} lists the Lean
  files and declarations of all numbered statements.\par\smallskip}

% (for ThinMatching/: \input{../scripts/overview_preamble.tex}) and is compiled from its own folder by
% scripts/build_overviews.py.
\documentclass[10pt,a4paper]{article}
\usepackage[T1]{fontenc}
\usepackage[utf8]{inputenc}
\usepackage[margin=2.2cm]{geometry}
\usepackage{amsmath,amssymb,amsthm}
\usepackage{array,booktabs,longtable}
\usepackage[dvipsnames]{xcolor}
\usepackage{enumitem}
\usepackage{tikz}
\usetikzlibrary{arrows.meta,positioning,shapes.multipart,fit,calc}
\usepackage[hidelinks]{hyperref}
\setlist{itemsep=2pt,topsep=4pt}
\renewcommand{\arraystretch}{1.2}

% Lean names, files and paper references
\newcommand{\lean}[1]{\textcolor{blue!50!black}{\texttt{\detokenize{#1}}}}
\newcommand{\file}[1]{\textcolor{OliveGreen!70!black}{\texttt{\detokenize{#1}}}}
\newcommand{\paperref}[1]{\textsf{#1}}
\newcommand{\hyp}[1]{\textcolor{BrickRed}{\textbf{#1}}}

% Boxes for the dependency diagrams
\tikzset{
  res/.style={draw, rounded corners, fill=blue!6, align=center, inner sep=4pt, font=\small},
  main/.style={res, fill=blue!18, very thick},
  cited/.style={res, fill=red!10, dashed},
  shared/.style={res, fill=gray!12},
  dep/.style={-{Stealth[length=2.2mm]}, thick},
}

\newcommand{\overviewtitle}[2]{%
  \begin{center}{\Large\bfseries #1}\\[4pt]{\small Proof overview of the Lean~4/Mathlib formalisation}\\[2pt]
  {\small #2}\end{center}}

% Standard paragraph describing the paper
\newcommand{\paperversion}{\noindent\textbf{Paper.} J.~Breitung, T.~Gervens, M.~Grohe.
  \emph{Graph Similarity Meets Optimal Transport}.\\ File \file{Graph_Similarity_Meets_Optimal_Transport.pdf} in the
  repository; statement, equation, page and line numbers refer to this PDF. \file{docs/NUMBERING.md} lists the Lean
  files and declarations of all numbered statements.\par\smallskip}

% (for ThinMatching/: \input{../scripts/overview_preamble.tex}) and is compiled from its own folder by
% scripts/build_overviews.py.
\documentclass[10pt,a4paper]{article}
\usepackage[T1]{fontenc}
\usepackage[utf8]{inputenc}
\usepackage[margin=2.2cm]{geometry}
\usepackage{amsmath,amssymb,amsthm}
\usepackage{array,booktabs,longtable}
\usepackage[dvipsnames]{xcolor}
\usepackage{enumitem}
\usepackage{tikz}
\usetikzlibrary{arrows.meta,positioning,shapes.multipart,fit,calc}
\usepackage[hidelinks]{hyperref}
\setlist{itemsep=2pt,topsep=4pt}
\renewcommand{\arraystretch}{1.2}

% Lean names, files and paper references
\newcommand{\lean}[1]{\textcolor{blue!50!black}{\texttt{\detokenize{#1}}}}
\newcommand{\file}[1]{\textcolor{OliveGreen!70!black}{\texttt{\detokenize{#1}}}}
\newcommand{\paperref}[1]{\textsf{#1}}
\newcommand{\hyp}[1]{\textcolor{BrickRed}{\textbf{#1}}}

% Boxes for the dependency diagrams
\tikzset{
  res/.style={draw, rounded corners, fill=blue!6, align=center, inner sep=4pt, font=\small},
  main/.style={res, fill=blue!18, very thick},
  cited/.style={res, fill=red!10, dashed},
  shared/.style={res, fill=gray!12},
  dep/.style={-{Stealth[length=2.2mm]}, thick},
}

\newcommand{\overviewtitle}[2]{%
  \begin{center}{\Large\bfseries #1}\\[4pt]{\small Proof overview of the Lean~4/Mathlib formalisation}\\[2pt]
  {\small #2}\end{center}}

% Standard paragraph describing the paper
\newcommand{\paperversion}{\noindent\textbf{Paper.} J.~Breitung, T.~Gervens, M.~Grohe.
  \emph{Graph Similarity Meets Optimal Transport}.\\ File \file{Graph_Similarity_Meets_Optimal_Transport.pdf} in the
  repository; statement, equation, page and line numbers refer to this PDF. \file{docs/NUMBERING.md} lists the Lean
  files and declarations of all numbered statements.\par\smallskip}

\begin{document}
\overviewtitle{Theorem 11: upper bound for the cut distance}{Folder \file{GsimOT/Theorem13_CutUpperBound}}
\paperversion

\section{Statement}
\textbf{Paper} (Theorem~11, p.~10, l.~333--337; proof overview pp.~10--11, l.~338--380; proof in App.~E.2, pp.~32--44).
There is a universal constant $C>0$ such that for all graphs $G,H$ of the same order, with
$\delta_\square=\delta_\square(G,H)$ and $\delta^{\mathrm{OT}}_\square=\delta^{\mathrm{OT}}_\square(G,H)$,
\[
\delta^{\mathrm{OT}}_\square\le\delta_\square\le C\,\delta^{\mathrm{OT}}_\square\ln\frac{e}{\delta^{\mathrm{OT}}_\square}
\ln\Bigl(e\ln\frac{e}{\delta^{\mathrm{OT}}_\square}\Bigr),
\]
the right-hand side being $0$ if $\delta^{\mathrm{OT}}_\square=0$. Consequently $\delta_\square=O((\delta^{\mathrm{OT}}_\square)^{1-\varepsilon})$
for every $\varepsilon>0$; Theorem~2 (p.~3, l.~118--120) states this as $\delta_\square\le C_\varepsilon(\delta^{\mathrm{OT}}_\square)^{1-\varepsilon}$,
and E.2.3 (p.~44, l.~1377--1379) as $\delta_\square\le C_\alpha(\delta^{\mathrm{OT}}_\square)^\alpha$ for every $\alpha<1$.

\textbf{Lean.} The paper cites the thin-matching theorem of Haqi and Oveis Gharan (Theorem~35 = [HOG26, Thm.~1.4]); it
is the hypothesis \lean{hThin : ThinPermutationTheorem}.
\begin{itemize}
\item \lean{GsimOT.theorem13_of_thin hThin}: $\exists C>0$ such that for all finite types $V,W$, graphs $G$ on $V$, $H$ on $W$
with $|V|=|W|$: $\delta^{\mathrm{OT}}_\square\le\delta_\square$ and $\delta_\square\le C f\ln(e/f)\ln(e\ln(e/f))$ with
$f=\delta^{\mathrm{OT}}_\square$. (In Lean $\ln 0=0$ and $e/0=0$, so the right-hand side is $0$ for $f=0$, as in the paper.)
\item \lean{GsimOT.theorem13_power_lt_one_of_thin hThin}: for every $\alpha<1$ there is $C_\alpha>0$ with
$\delta_\square\le C_\alpha f^\alpha$ (the statement of E.2.3; \lean{theorem13_power_of_thin} is the special case $0<\alpha<1$).
\item \lean{GsimOT.theorem2_cut_of_thin hThin}: for every $\varepsilon>0$ there is $C_\varepsilon>0$ with
$\delta_\square\le C_\varepsilon f^{1-\varepsilon}$ (Theorem~2, cut part).
\item \lean{theorem13_of_rackeShah}, \lean{theorem13_power_lt_one_of_rackeShah}, \lean{theorem2_cut_of_rackeShah}
(\file{FromThinMatching.lean}): the same, assuming only the tree-cut sparsifier theorem [RS14], since Theorem~35 is
derived from the separate library \file{ThinMatching/}.
\end{itemize}
These are exactly the paper's statements (with the paper's constant $C=\max\{2,39C_1,113\}$, $C_1=1+4C_{\mathrm{thin}}$).

\section{Where the proof lives}
\begin{longtable}{>{\raggedright\arraybackslash}p{5.2cm} >{\raggedright\arraybackslash}p{10.6cm}}
\toprule \textbf{File} & \textbf{Content} \\ \midrule\endhead \bottomrule\endfoot
\file{Definitions.lean} & thinness Eq.~(23) (\lean{IsThin}), Theorem~35 as a hypothesis (\lean{ThinPermutationTheorem}),
the contraction principle statement, support functions, $d_n$ \\
\file{Theorem52_ThinRounding.lean} & App.~E.2.1: Lemma~36 (\lean{lemma50}), Lemma~37 (\lean{lemma51}), Eqs.~(24)--(28),
Eq.~(26) (\lean{theorem52}), Theorem~38 (\lean{theorem38}) \\
\file{Lemma46_SpectralTruncation.lean}, \file{Lemma47_MergingPermutations.lean}, \file{ContractionPrinciple.lean},
\file{Theorem45_AdditiveRounding.lean} & App.~E.2.2: the additive rounding, Theorem~39 (\lean{theorem45}) \\
\file{Theorem13.lean} & App.~E.2.3: graph facts, numerics, Eq.~(39) (\lean{eq33}), Theorem~11 (\lean{theorem13}),
the power bounds and Theorem~2 \\
\file{FromThinMatching.lean} & Theorem~35 from \file{ThinMatching/} and the \lean{_of_rackeShah} versions \\
\end{longtable}
The paper's $\zeta_Q$ (Eq.~(21)) is \lean{epsR (adjMat G) (adjMat H) Q}, and $\zeta_Q=n^{-2}C_\square(Q)$ (\lean{epsR_adjMat}).

\section{Proof layout}
\begin{center}
\begin{tikzpicture}[every node/.append style={font=\footnotesize}]
  \node[main] (rs) at (-0.8,0.2) {Theorem 11 / Theorem 2, assuming [RS14]\\\lean{theorem13_of_rackeShah}, \lean{theorem2_cut_of_rackeShah}};
  \node[main] (t11) at (-0.8,-1.4) {Theorem 11 (E.2.3)\\\lean{theorem13}, \lean{theorem13_power_lt_one}};
  \node[res] (num) at (5.6,-1.4) {numerics of E.2.3\\\lean{main_numeric}, \lean{power_numeric_lt_one}};
  \node[res] (e39) at (-3.6,-3.0) {Eq. (39), $C_1=1+4C_{\mathrm{thin}}$\\\lean{eq33}};
  \node[res] (t39) at (3.6,-3.0) {Theorem 39 (additive)\\\lean{theorem45}};
  \node[res] (t38) at (-3.6,-4.5) {Theorem 38\\\lean{theorem38}};
  \node[res] (e26) at (-3.6,-6.0) {Eq. (26): $\zeta_P\le(1+4\alpha)\zeta_Q$\\\lean{theorem52}};
  \node[res] (e27l) at (-6.4,-7.6) {Eq. (27), first\\\lean{eq27_left}};
  \node[res] (e27r) at (-2.4,-7.6) {Eq. (27), second; Eq. (28)\\\lean{eq27_right}, \lean{eq28_cut}};
  \node[res] (ea) at (-4.6,-9.1) {$\|EA\|_\square\le2\alpha n^2\zeta_Q$, Eq. (25)\\\lean{normEA_le}, \lean{cut_R_EAt_le}};
  \node[res] (l36) at (-6.4,-10.6) {Lemma 36\\\lean{lemma50}};
  \node[res] (l37) at (-3.0,-10.6) {Lemma 37, Eq. (24)\\\lean{lemma51}};
  \node[cited] (hog) at (-7.0,-4.5) {Theorem 35 [HOG26]\\\lean{ThinPermutationTheorem}};
  \node[res] (l48) at (3.6,-4.5) {Lemma 48 (tree merging)\\\lean{lemma48}};
  \node[res] (l47) at (3.6,-6.0) {Lemma 47 (two permutations)\\\file{Lemma47_MergingPermutations.lean}};
  \node[res] (c46) at (6.9,-8.6) {Lemma 45, Corollary 46\\\file{Lemma46_SpectralTruncation.lean}};
  \node[res] (cp) at (2.4,-7.6) {Lemmas 42--44, [LT91]\\\lean{contractionPrinciple}};
  \node[shared] (com) at (2.0,-10.6) {Common: cut norm, $\mathrm{DS}$, Birkhoff\\\lean{cutNorm}, \lean{epsR}, \lean{IsDS}};
  \node[cited] (rsc) at (6.0,0.2) {[RS14] via \file{ThinMatching/}\\\lean{RackeShahTheorem}};
  \draw[dep] (rs) -- (t11); \draw[dep] (rs) -- (rsc);
  \draw[dep] (t11) -- (num); \draw[dep] (t11) -- (e39); \draw[dep] (t11) -- (t39);
  \draw[dep] (e39) -- (t38); \draw[dep] (t38) -- (e26); \draw[dep] (t38) -- (hog);
  \draw[dep] (e26) -- (e27l); \draw[dep] (e26) -- (e27r); \draw[dep] (e27r) -- (ea);
  \draw[dep] (ea) -- (l36); \draw[dep] (ea) -- (l37);
  \draw[dep] (t39) -- (l48); \draw[dep] (l48) -- (l47); \draw[dep] (t39.east) to[bend left=30] (c46);
  \draw[dep] (l47) -- (cp);
  \draw[dep] (l37) -- (com); \draw[dep] (cp) -- (com); \draw[dep] (c46) -- (com);
\end{tikzpicture}
\end{center}
Left: the multiplicative branch (App.~E.2.1); right: the additive branch (App.~E.2.2); top: the combination (App.~E.2.3).
Dashed boxes are cited results, which are explicit hypotheses.

\section{The multiplicative bound (App. E.2.1, pp. 32--35)}
Throughout App.~E.2, $V=W=[n]$ with $n\ge2$, $A,B$ are the adjacency matrices, and $\zeta_Q$ is defined by Eq.~(21)
(sets) $=$ Eq.~(22) ($[0,1]$-valued functions), p.~32. In Lean $V$, $W$ are arbitrary finite types of equal size;
testing on sets is \lean{abs_cut_le_epsR}.
\begin{enumerate}
\item \textbf{Thinness, Eq. (23)} (p.~32, l.~1014--1015): $P$ is $\alpha$-thin w.r.t.\ $Q$ if for all $X,Y$ the mass of $P$ on the cut
$(X\times\overline Y)\cup(\overline X\times Y)$ is at most $\alpha$ times that of $Q$. \lean{IsThin} is this, written with
$U=U_V\sqcup U_W$ as in the proof overview (p.~10, l.~350--353).
\item \textbf{Theorem 35} (p.~32, l.~1021--1022) is the hypothesis \lean{ThinPermutationTheorem}: some $P$ is
$C_{\mathrm{thin}}\ln(2n)\ln\ln(4n)$-thin.
\item \textbf{Lemma 36} (p.~33, l.~1026--1028): $\sum_i|x_i-y_{\pi(i)}|\le\alpha\sum_{ij}Q_{ij}|x_i-y_j|$. As in the paper: Eq.~(23) for the
level sets $X_s,Y_s$ gives the pointwise inequality (\lean{hpt} in \lean{lemma50}, with \lean{levelInd_apply}), and
$\int|1[x>s]-1[y>s]|\,ds=|x-y|$ (\lean{integral_levelInd}); integrate.
\item \textbf{Lemma 37} (p.~33, l.~1038--1039; proof pp.~33--34): $n^{-2}\sum_{ij}Q_{ij}|x_i-(Q^\top x)_j|\le4\zeta_Q$ for $x=A\mathbf 1_T$
(\lean{lemma51}; the paper allows $y\in[0,1]^n$, only $y=\mathbf1_Y$ is used). Steps: $F_{ij}=x_i-r_j$, Eq.~(24)
$n^{-2}\sum Q_{ij}|F_{ij}|\le2\zeta_Q$ via $S_\pm$, then $x_i-(Q^\top x)_j=F_{ij}-\sum_{i'}Q_{i'j}F_{i'j}$ and the triangle
inequality.
\item \textbf{Eq. (25)} (p.~34, l.~1066--1072): $\|RK\|_\square\le\|K\|_\square$, used in the form \lean{abs_sum_mul_le_of_cut}.
\item \textbf{Theorem 38} (p.~34, l.~1074--1075; proof pp.~34--35): $\zeta_P\le(1+4\alpha)\zeta_Q$ for some $P$ (\lean{theorem38}). Its proof:
  \begin{enumerate}
  \item $\zeta_P=n^{-2}\|A-PBP^\top\|_\square$ (p.~34, l.~1080--1084): \lean{epsR_permMat_le} (direction $\le$).
  \item split $A-PBP^\top=(A-RAR^\top)+P(Q^\top AQ-B)P^\top$, $R=PQ^\top$, i.e.\ $R(v,k)=Q(k,\pi v)$ (p.~35, l.~1085--1089): \lean{hsplit} in
  \lean{theorem52}; conjugation invariance by testing on $\pi(S)\times\pi(T)$.
  \item Eq.~(27), first bound $n^{-2}\|Q^\top AQ-B\|_\square\le\zeta_Q$ (p.~35, l.~1091--1095): \lean{eq27_left}.
  \item $f=EA\mathbf1_Y$ with $E=I-R$; Lemma~36, Lemma~37, $\sum_if_i=0$ give $n^{-2}\|EA\|_\square\le2\alpha\zeta_Q$
  (l.~1096--1105): \lean{normEA_le}, with \lean{abs_sum_le_half_of_sum_eq_zero}.
  \item Eq.~(28) $A-RAR^\top=EA+R(EA)^\top$ (symmetry of $A$): \lean{eq28_cut}; Eq.~(25) and transposition:
  \lean{cut_R_EAt_le}; together the second bound of Eq.~(27): \lean{eq27_right}.
  \item Eq.~(26) for every $\alpha$-thin $P$: \lean{theorem52}; with Theorem~35: \lean{theorem38}.
  \end{enumerate}
\end{enumerate}

\IfFileExists{Theorem13_overview_additive.tex}{% Fragment: the additive part of Theorem 11 (Theorem 39, Appendix E.2.2).
% It is \input by Theorem13_overview.tex; no preamble here.
\section{The additive bound (Theorem 39)}

\subsection{Statement}

\paragraph{Paper (Theorem~39, p.~36, l.~1112--1114, App.~E.2.2).} Let $d_n:=\lceil 2\log n\rceil$
($\log$ is the binary logarithm, $n\ge2$). For every $Q\in\mathrm{DS}(n)$ there is a permutation matrix $P$ with
\[ \zeta_P\le\zeta_Q+36\sqrt{d_n}\,n^{-1/6}. \]

\paragraph{Lean.} \lean{GsimOT.theorem45} (file \file{Theorem45_AdditiveRounding.lean}):
for finite types $V,W$ with $|V|=|W|=n\ge2$, symmetric matrices $A\in\mathbb R^{V\times V}$, $B\in\mathbb R^{W\times W}$
with entries in $[0,1]$ and $Q\in\mathrm{DS}(V,W)$ there is a bijection $\pi\colon V\to W$ with
\[ \lean{epsR A B (permMat pi)}\le\lean{epsR A B Q}+36\sqrt{d_n}\,n^{-1/6},\qquad d_n=\lean{dn n}=\lceil\log_2(n^2)\rceil. \]
Here \lean{epsR A B R} is $\zeta_R(A,B)$, the weighted cut norm of the discrepancy kernel on $(\Omega_R,\mu_R)$;
for adjacency matrices it is $\zeta_R$ (Eq.~(21)). The hypothesis \hyp{hCP : ContractionPrinciple} (Lemma~43,
[LT91]) is discharged by \lean{GsimOT.contractionPrinciple}. The statement is the paper's (for adjacency matrices
of graphs of the same order), slightly more general (any symmetric $[0,1]$-matrices).

\subsection{Files}
\begin{longtable}{>{\raggedright\arraybackslash}p{6.3cm} >{\raggedright\arraybackslash}p{9.5cm}}
\toprule \textbf{File} & \textbf{Content} \\ \midrule\endhead \bottomrule\endfoot
\file{Definitions.lean} & \lean{radAvg}, \lean{ContractionPrinciple} (statement of Lemma~43), \lean{hSupp} ($h_Q$, Eq.~(33)), \lean{rhoSupp} ($\rho$, Eq.~(34)), \lean{dn} \\
\file{ContractionPrinciple.lean} & Lemma~43 proved: \lean{contractionPrinciple} (Talagrand's comparison argument) \\
\file{Lemma46_SpectralTruncation.lean} & Lemma~40, spectral truncation and Eq.~(31), Lemma~45 (Eq.~(30)), zonotopes and Eqs.~(32), (33), (36), (37), Lemma~41, Corollary~46 \\
\file{Lemma47_MergingPermutations.lean} & Lemma~42 (\lean{Lemma47.lemma43}), Lemma~44, the blocks, Eq.~(38), Lemma~47 \\
\file{Theorem45_AdditiveRounding.lean} & Eq.~(35), Eq.~(3) with Carath\'eodory, Lemma~48, Theorem~39 \\
\end{longtable}

\subsection{Proof layout}
\begin{center}
\begin{tikzpicture}[every node/.append style={font=\footnotesize}]
  \node[main] (t39) at (0,0) {Theorem 39\\\lean{theorem45}};
  \node[res] (l48) at (-4,-1.7) {Lemma 48\\\lean{lemma48}};
  \node[res] (c46) at (3.2,-1.7) {Corollary 46\\\lean{corollary46}};
  \node[res] (car) at (-6.4,-3.4) {Eq.~(3) + Carath\'eodory\\\lean{exists_perm_decomp}};
  \node[res] (l47) at (-2.6,-3.4) {Lemma 47\\\lean{lemma47}};
  \node[res] (l41) at (0.5,-3.4) {Lemma 41\\\lean{lemma41}};
  \node[res] (z) at (3.4,-3.4) {Eqs.~(33), (36), (37)\\\lean{eq36}, \lean{eq37}};
  \node[res] (l45) at (7.0,-3.4) {Lemma 45, Eq.~(30)\\\lean{lemma45}};
  \node[res] (l42) at (-4.5,-5.1) {Lemma 42\\\lean{Lemma47.lemma43}};
  \node[res] (l44) at (-1.2,-5.1) {Lemma 44\\\lean{Lemma47.lemma44}};
  \node[res] (l40) at (7.0,-5.1) {Lemma 40, Eq.~(31)\\\lean{lemma41_2_fst}, \lean{eq28}};
  \node[res] (l43) at (-1.2,-6.8) {Lemma 43 [LT91]\\\lean{contractionPrinciple}};
  \node[shared] (cut) at (7.0,-6.8) {cut norm (Common)\\\lean{cutNorm}, \lean{epsR}};
  \draw[dep] (t39) -- (l48); \draw[dep] (t39) -- (c46);
  \draw[dep] (t39) to[bend left=22] (l45);
  \draw[dep] (l48) -- (l47); \draw[dep] (l48) -- (car);
  \draw[dep] (l47) -- (l42); \draw[dep] (l47) -- (l44); \draw[dep] (l44) -- (l43);
  \draw[dep] (c46) -- (l41); \draw[dep] (c46) -- (z); \draw[dep] (c46) -- (l45);
  \draw[dep] (l45) -- (l40); \draw[dep] (l40) -- (cut);
\end{tikzpicture}
\end{center}
Lemma~48 also uses Eq.~(35) (\lean{rhoSupp_convex}); all results in this part are proved, the only external input
(Lemma~43, cited in the paper) is proved in \file{ContractionPrinciple.lean}.

\subsection{The proof step by step}

\paragraph{Notation (p.~36).} The paper works in $\mathcal L^2(n)$ (inner product $\frac1n\sum_ix_iy_i$) with
averaging operators. Lean uses the standard inner product and rescales: the paper's eigenvalues are $\lambda_r/n$,
its eigenfunctions $\sqrt n\,e_r$, and $\|K\|_{\mathcal L^2(n)}\le t$ becomes \lean{OpBound K (t * n)}
($|\sum_{i,k}f_iK_{ik}g_k|\le tn\cdot\frac{\|f\|^2+\|g\|^2}2$, equivalent by scaling). The couplings $\mu_Q(i,j)=Q_{ij}/n$
and $\zeta_Q(A',B')=\|D_{A',B'}\|_{\square,\mu_Q}$ are built into \lean{epsR}.

\paragraph{Lemma 40 (Lifted Kernels, p.~36).} \lean{lemma41_2_fst}, \lean{lemma41_2_snd}. As in the paper, for
$S,T\subseteq\Omega_Q$ the marginal weights $s_i=\sum_jQ_{ij}\mathbf 1_S(i,j)\in[0,1]$ have $\|s\|\le1$
(\lean{weight_sq_le}), and the cut value is $\langle s,K\sigma\rangle$, bounded by the operator norm.

\paragraph{Lemma 41 (Close Support Functions, p.~37).} \lean{lemma41}, for non-empty compact
$Z,Z'$ in any real inner product space (without the convexity hypothesis), $\|T\|\le\tau$, with $h_Z$ = \lean{supportFn} and
$\max_{u,u'\in Z}|\langle u,Tu'\rangle|$ = \lean{bilMax}. The paper's steps: $|h_Z(w)-h_{Z'}(w)|\le\eta\|w\|$
(\lean{abs_supportFn_sub_le}, by homogeneity \lean{supportFn_smul}); $m_T(Z,Z)\le m_T(Z',Z)+M\eta\tau$ via
$h_Z(Tu')$, and $m_T(Z',Z)\le m_T(Z',Z')+M\eta\tau$ via $h_{Z'}(T^*u)$ with the adjoint (\lean{lemma41_step});
the same for $-T$ gives the absolute values (\lean{abs_inner_le_bilMax_add}).

\paragraph{Lemma 42 (symmetrisation, pp.~37--38).} \lean{Lemma47.lemma43}: independent $X_a\in\{0,1\}$ with
$\mathbb E X_a=\theta_a$ (product weights \lean{bern}), bounded $d_a$; independent copy $X'$ and Jensen, symmetry of
$X_a-X'_a$ (coordinate swap \lean{swapAt}), triangle inequality --- as in the paper.

\paragraph{Lemma 43 (Talagrand Contraction Lemma, p.~38).} Cited in the paper; stated as \lean{ContractionPrinciple}
(for non-empty bounded $S$) and proved by \lean{contractionPrinciple}: the two-point lemma, the comparison
theorem without absolute values (one coordinate at a time), then absolute values and the factor~$2$ by adding $0$
to $S$ and using $\sigma\mapsto-\sigma$.

\paragraph{Lemma 44 (p.~38, proof pp.~38--39, Eq.~(29)).} \lean{Lemma47.lemma44}: for non-empty blocks $B_a$ and
$\|u_i\|\le M$, $\mathbb E_\sigma\sup_{w\in B_H}|\sum_a\sigma_a\sum_{i\in B_a}\langle w,u_i\rangle_+|\le2M(\sum_a\ell_a^2)^{1/2}$.
As in the paper: $\sum_{i\in B_a}\langle w,u_i\rangle_+=\mathbb E_{I_a}(\ell_a\langle w,u_{I_a}\rangle)_+$ with uniform
independent $I_a\in B_a$ (\lean{sum_piFinset_comp}, over \lean{Fintype.piFinset B}); Jensen; for fixed $I$ Lemma~43 with
$\gamma(s)=s_+$; $\sup_w|\langle w,g_\sigma\rangle|\le\|g_\sigma\|$; $\mathbb E\|g_\sigma\|\le(\mathbb E\|g_\sigma\|^2)^{1/2}$
(\lean{radAvg_norm_le}).

\paragraph{Lemma 45 (spectral approximation, p.~39, proof pp.~39--40, Eqs.~(30), (31)).} \lean{lemma45}. Spectral decomposition and
Parseval (\lean{spectral_entry}, \lean{parseval_row}); truncations $A_t,B_t$ (\lean{specTrunc}) with
$\|A-A_t\|_{\mathcal L^2(n)}\le t$ (\lean{opBound_sub_specTrunc}); Eq.~(31), $|\zeta_Q(A,B)-\zeta_Q(A_t,B_t)|\le2t$, from
$D_{A,B}=D_{A_t,B_t}+(A-A_t)^{(1)}-(B-B_t)^{(2)}$ and Lemma~40 (\lean{eq28}); the vectors $z_{ij}$ (\lean{specVec},
$\|z_{ij}\|\le\sqrt2$ by Parseval, \lean{norm_specVec_le}) and $T_t=\mathrm{diag}(\lambda_r^{-1},-\nu_s^{-1})$
(\lean{diagOp}, \lean{specDiag}, $\|T_t\|\le1/t$ by \lean{norm_diagOp_le}) with
$\langle z_{ij},T_tz_{k\ell}\rangle=(A_t)_{ik}-(B_t)_{j\ell}$ (\lean{diagForm_specVec}, \lean{inner_diagOp}).
The maximum in Eq.~(30) is the cut norm of \lean{zonKernel}. Lean takes $H=\mathbb R^{V\sqcup W}$ (the coordinates
outside $I_A\sqcup I_B$ are zero), so the case $I_A=I_B=\emptyset$ is not separate.

\paragraph{Zonotopes and $\rho$ (pp.~40--41, Eqs.~(32)--(37)).} $Z_Q$ is the image of \lean{cube} under \lean{zonLin v Q} (Eq.~(32)),
compact, convex, in the ball of radius $M$ (\lean{isCompact_zonotope}, \lean{convex_zonotope},
\lean{norm_zonLin_le}); Eq.~(33), $h_{Z_Q}=h_Q$ (\lean{supportFn_zonotope}, maximiser $a_{ij}=\mathbf 1[\langle w,z_{ij}\rangle>0]$);
Eq.~(34) is \lean{rhoSupp}; Eq.~(35) (\lean{rhoSupp_triangle}, \lean{rhoSupp_convex}); Eq.~(36) as an identity
(\lean{eq36}: bilinear expansion \lean{eq36_expand}, maximum at indicator vectors); Eq.~(37) (\lean{eq37}).

\paragraph{Corollary 46 (p.~41).} \lean{corollary46}: for $v,T$ satisfying Eq.~(30) and all $P,Q\in\mathrm{DS}(n)$,
$\zeta_P\le\zeta_Q+4t+\frac{2\sqrt2}{t}\rho(P,Q)$: Lemma~41 with $M=\sqrt2$, $\eta=\rho(P,Q)$, $\tau=1/t$, and Eq.~(37)
for $P$ and $Q$ --- exactly the paper's chain.

\paragraph{Lemma 47 (Merging Two Permutations, pp.~41--43).} \lean{lemma47} (statement, $20Mn^{-1/3}$) from
\lean{lemma47_sharp} ($18Mn^{-1/3}$, what the proof gives on p.~43, l.~1333--1334). Steps: $\tau=\pi_0^{-1}\pi_1$ and
$u_i=z_{i,\pi_0(i)}$ (the w.l.o.g.\ $P_0=I$); $L=\lceil n^{1/3}\rceil$; blocks with properties (i), (ii)
(\lean{exists_blocks}, from a maximal $L$-separated set of cut points, \lean{exists_sep_cover}); property (i) gives
$\sum_a|B_a|^2\le2Ln$; $n(h_{T_X}-(1-\theta)h_{\mathrm{id}}-\theta h_\tau)=\sum_a(X_a-\theta)d_a$; Lemma~42 and twice
Lemma~44 give some $X$ with $\Delta(T_X)\le8M\sqrt{2L/n}\le16Mn^{-1/3}$ (Eq.~(38), \lean{exists_good_choice}); repair
of $T_X$ outside the at most $n/L$ block boundaries (\lean{exists_equiv_eq_off}), each changed row costing
$\le2M/n$ (\lean{phi_sub_le}).

\paragraph{Lemma 48 (p.~43).} \lean{lemma48}: Birkhoff (Eq.~(3)) and Carath\'eodory give
$Q=\sum_{s\le N}\lambda_sP_s$ with $\lambda_s>0$, $N\le n^2$ (\lean{exists_perm_decomp}, via Mathlib's Birkhoff theorem and
\lean{eq_pos_convex_span_of_mem_convexHull}; the bound $N\le n^2$ from affine independence in $\mathbb R^{n\times n}$
with constant total mass); a binary tree of height $\le d_n$ (\lean{merge_aux}: induction on the height $k$, splitting the
index set in halves; \lean{combo} is $Q_\nu$); at each node Lemma~47, the triangle inequality and Eq.~(35) give
$\rho(P_\nu,Q_\nu)\le\eta_0+(k-1)\eta_0$.

\paragraph{Proof of Theorem 39 (p.~43).} \lean{theorem45}: $t=\sqrt{20d_n}\,n^{-1/6}$, Lemma~45 for this $t$,
Lemma~48 ($\rho(P,Q)\le20\sqrt2\,d_nn^{-1/3}$), Corollary~46: $\zeta_P\le\zeta_Q+4t+80d_nn^{-1/3}/t=\zeta_Q+8t$ and
$8\sqrt{20}<36$ (\lean{hs20}: $\sqrt{20}\le9/2$).

\subsection{Notes on the formalisation (additive part)}
\begin{itemize}
\item Standard inner product instead of $\mathcal L^2(n)$; $H=\mathbb R^{V\sqcup W}$ instead of $\mathbb R^{I_A\sqcup I_B}$.
\item Lemma~41 is stated without the convexity hypothesis and Lemma~44 without the disjointness hypothesis;
  the paper's arguments prove these more general forms.
\item Lemma~47: blocks from a maximal separated set instead of the cycle decomposition (only properties (i), (ii)
  are used); Lemma~44 is applied to all blocks with vectors $X_au_i$ instead of only the blocks with $X_a=1$;
  the proof gives the constant $18M$, as in the paper's proof.
\item Lemma~48: the binary tree halves the index set at every node instead of padding with zero-weight leaves.
\item Expectations are finite weighted sums (\lean{radAvg}, \lean{bern}, uniform averages over a \lean{piFinset}).
\end{itemize}
The pseudometric $\rho$ is formalised by its definition via support functions, Eq.~(34) (\lean{rhoSupp}); for the
compact convex zonotopes $Z_Q$ this is the Hausdorff distance of the proof overview (p.~11, l.~368--369).

\subsection{How to check}
\begin{itemize}
\item \texttt{lake build GsimOT.Theorem13\_CutUpperBound.Theorem45\_AdditiveRounding}
\item \texttt{\#print axioms GsimOT.theorem45}: only \lean{propext}, \lean{Classical.choice}, \lean{Quot.sound}.
\end{itemize}
}{%
\section{The additive bound (App. E.2.2)}
Theorem~39 (p.~36): $\zeta_P\le\zeta_Q+36\sqrt{d_n}\,n^{-1/6}$, $d_n=\lceil2\log n\rceil$, is \lean{theorem45}
(\file{Theorem45_AdditiveRounding.lean}), assuming the contraction principle, which is proved in
\file{ContractionPrinciple.lean}. (A detailed overview of this part is written separately.)}

\section{Combining both bounds: the proof of Theorem 11 (App. E.2.3, p. 44)}
\begin{enumerate}
\item \emph{Lower bound} (l.~1363): permutation matrices are doubly stochastic: \lean{deltaOTCut_le_deltaCut}.
\item \emph{Optimal $Q$}: ``let $Q\in\mathrm{DS}(n)$ be optimal, so that $\zeta_Q=\varepsilon$'': attained by compactness of
$\mathrm{DS}(V,W)$ and continuity of $C_\square$ (\lean{isCompact_DS}, \lean{continuous_costCut},
\lean{exists_costCut_eq_dOTCut}).
\item \emph{Eq. (39)} (l.~1364--1366): $\ln(2n)\ln\ln(4n)\ge\ln4\cdot\ln\ln8>1$ for $n\ge2$
(\lean{one_lt_log_four_mul_log_log_eight}, \lean{one_le_log_mul_log_log}), so Theorem~38 gives
$\delta_\square\le C_1\ln(2n)\ln\ln(4n)\,\varepsilon$ with $C_1=1+4C_{\mathrm{thin}}$: \lean{eq33}.
\item \emph{Cases} (all in \lean{main_numeric}): $\varepsilon=0$ (from Eq.~(39)); $\varepsilon\ge\frac12$ ($\delta_\square\le1$, $C\ge2$);
$0<\varepsilon<\frac12$ with $L_\varepsilon=\ln(e/\varepsilon)$:
  \begin{itemize}
  \item $n\le\varepsilon^{-12}$: $\ln(2n)\le13L_\varepsilon$ and $\ln\ln(4n)\le\ln(14L_\varepsilon)\le3\ln(eL_\varepsilon)$ using $\ln14<3$
  (\lean{log_fourteen_lt_three}), hence $C=39C_1$;
  \item $n>\varepsilon^{-12}$: $d_n\le3\log n$ (\lean{dn_mul_log_two_le}) and the error term of Theorem~39 is $<112\varepsilon$
  (\lean{error_term_le}), hence $\delta_\square\le113\varepsilon$.
  \end{itemize}
  So $C=\max\{2,39C_1,113\}$ (\lean{theorem13}).
\item \emph{Power bound} (l.~1377--1379): $\varepsilon^{1-\alpha}\ln(e/\varepsilon)\ln(e\ln(e/\varepsilon))$ is bounded on $(0,1]$:
\lean{power_numeric} ($0<\alpha<1$, with constant $(1+\frac2{1-\alpha})^2$) and \lean{power_numeric_lt_one} (every $\alpha<1$;
for $\alpha\le\frac12$ via $\alpha=\frac12$ and $f^{1/2}\le f^\alpha$). Results: \lean{theorem13_power_lt_one},
\lean{theorem2_cut_of_thin} ($\alpha=1-\varepsilon$).
\end{enumerate}

\section{Assumptions and trusted base}
\begin{itemize}
\item \hyp{Theorem 35} [HOG26, Thm.~1.4] is the hypothesis \lean{ThinPermutationTheorem} of the \lean{_of_thin} theorems; it is
derived from [RS14] (\lean{RackeShahTheorem}, cited in [HOG26]) in the \file{ThinMatching/} library, so the
\lean{_of_rackeShah} theorems assume only [RS14].
\item The contraction principle [LT91] (Lemma~43) is a hypothesis of \lean{theorem13}, but it is proved
(\lean{contractionPrinciple}), so the \lean{_of_thin} theorems do not assume it.
\item Only the standard axioms \lean{propext}, \lean{Classical.choice}, \lean{Quot.sound} are used.
\end{itemize}

\section{Notes on the formalisation}
\begin{itemize}
\item Lemma~37 is proved for $y=\mathbf1_T$ only (the case used).
\item The Lean hypothesis gives thinness with exactly $\alpha=C_{\mathrm{thin}}\ln(2n)\ln\ln(4n)$, so the monotonicity step of the
proof of Theorem~38 (``$\alpha'$-thin, $\alpha'\le\alpha$, hence $\alpha$-thin'') is not needed.
\item $n=1$ (excluded by ``$n\ge2$'' in App.~E.2) is handled separately: $\delta_\square=0$
(\lean{deltaCut_le_zero_of_card_eq_one}, \lean{rhs_nonneg}).
\item For $n>\varepsilon^{-12}$ the paper uses that $x\mapsto x^{-1/6}\sqrt{\ln x}$ is decreasing and
$\varepsilon\sqrt{\ln(1/\varepsilon)}\le(2e)^{-1/2}$; Lean uses $m=n^{1/6}>\varepsilon^{-2}$ and $\ln m\le m/e$ (the same
inequality), with the same constant $112$.
\end{itemize}

\section{How to check this result}
\begin{verbatim}
lake build GsimOT.Theorem13_CutUpperBound.Theorem13 GsimOT.Theorem13_CutUpperBound.FromThinMatching
lake env lean scripts/CheckAxioms.lean   # axioms of theorem13_of_thin, theorem13_of_rackeShah, ...
\end{verbatim}
The correspondence table is \file{Theorem13_correspondence.pdf} in the same folder.
\end{document}
